% Generated draft: requires mathematical and source review.
\documentclass[11pt]{article}
\usepackage[margin=25mm]{geometry}
\usepackage{fontspec}
\setmainfont{Latin Modern Roman}
\usepackage{amsmath,amssymb,mathtools,unicode-math}
\setmathfont{Latin Modern Math}
\usepackage{enumitem,needspace,xurl,hyperref}
\setlength{\parindent}{0pt}
\setlength{\parskip}{5pt}
\setlength{\emergencystretch}{3em}
\newcommand{\N}{\mathbb{N}}
\newcommand{\Z}{\mathbb{Z}}
\newcommand{\Q}{\mathbb{Q}}
\newcommand{\R}{\mathbb{R}}
\newcommand{\C}{\mathbb{C}}
\newcommand{\F}{\mathbb{F}}
\newcommand{\A}{\mathbb{A}}
\newcommand{\PP}{\mathbb{P}}
\newcommand{\E}{\mathbb{E}}
\newcommand{\eps}{\varepsilon}
\newcommand{\dd}{\,\mathrm{d}}
\begin{document}
\section*{UnsolvedMath ID: 701638}
\textbf{Source problem code:} catalogue entry
\subsection{Definitions and mathematical statement}
A \emph{skew brace} is a set \(B\) equipped with binary operations \(+\) and \(\cdot\) such that \((B,+)\) and \((B,\cdot)\) are groups and
\[
 a\cdot(b+c)=(a\cdot b)+(-a)+(a\cdot c)
 \qquad\text{for all }a,b,c\in B,
\]
where \(-a\) denotes the inverse in \((B,+)\). The additive group is not assumed to be abelian, and the order of the three summands in this identity is part of the definition.

For a group \(G\), written multiplicatively for these definitions, put \([x,y]=x^{-1}y^{-1}xy\), and let \([H,K]\) denote the subgroup generated by all \([h,k]\) with \(h\in H\) and \(k\in K\). Define
\[
 G^{(0)}=G,\qquad G^{(n+1)}=[G^{(n)},G^{(n)}],
 \qquad
 \gamma_1(G)=G,\qquad \gamma_{n+1}(G)=[\gamma_n(G),G].
\]
The group \(G\) is \emph{soluble} if \(G^{(d)}=\{e\}\) for some integer \(d\geq0\), and is \emph{nilpotent} if \(\gamma_{c+1}(G)=\{e\}\) for some integer \(c\geq0\). Non-soluble means that no such finite derived length exists. These definitions apply to both groups of a skew brace.

The supplied record identifies the following as the remaining questions of Problem 19.90:
\begin{itemize}
\item[(b)] Does there exist a skew brace \((B,+,\cdot)\) such that \((B,+)\) is non-soluble and \((B,\cdot)\) is nilpotent? No finiteness restriction is imposed on \(B\).
\item[(c)] Does there exist a skew brace \((B,+,\cdot)\) with \(|B|<\infty\) such that \((B,+)\) is soluble and \((B,\cdot)\) is non-soluble?
\end{itemize}
The same record includes the following original questions as archived solved parts:
\begin{itemize}
\item[(a)] Does there exist a skew brace \((B,+,\cdot)\), without a finiteness restriction, such that \((B,+)\) is soluble and \((B,\cdot)\) is non-soluble?
\item[(d)] Does there exist a skew brace \((B,+,\cdot)\) with \(|B|<\infty\) such that \((B,+)\) is non-soluble and \((B,\cdot)\) is nilpotent?
\end{itemize}
The supplied archive reports an affirmative answer to (a), attributed to T. Nasybullov, and a negative answer to (d), attributed to C. Tsang and Q. Chao. These are supplied source assessments, not independently verified conclusions of this draft. The supplied overall status is \texttt{partially\_solved}.
\subsection{Short English statement}
The remaining questions ask whether a skew brace can have a non-soluble additive group and a nilpotent multiplicative group, with no finiteness restriction in (b), and whether a finite skew brace can have a soluble additive group and a non-soluble multiplicative group in (c). The supplied archive separately reports that (a) has an affirmative answer and (d) a negative answer; those reports have not been independently verified.
\subsection{Sources}
\begin{itemize}
\item Local UnsolvedMath registry, supplied complete record: permanent ID \(701638\), exact source problem code \texttt{catalogue entry}, title \emph{Kourovka Notebook Problem 19.90}. Proposed by A. Smoktunowicz and L. Vendramin in 2018.
\item \emph{The Kourovka Notebook (Unsolved Problems in Group Theory)}, 19th issue (2018), Problem 19.90; the supplied background states that it was updated in the 21st issue (2026), identified as \texttt{arXiv:1401.0300v48}. The registry describes the collection as Issues 1--21 (1965--2026), edited by E. I. Khukhro and V. D. Mazurov.
\item Supplied Archive of Solved Problems attribution for (a): T. Nasybullov, \emph{J. Algebra} \textbf{540} (2019), 156--167; reported answer: yes. No article title, DOI, or article URL is supplied.
\item Supplied Archive of Solved Problems attribution for (d): C. Tsang and Q. Chao, \emph{Int. J. Algebra Comput.} \textbf{30}, no. 2 (2020), 253--265; reported answer: no. No article title, DOI, or article URL is supplied.
\item Supplied status review checked on \texttt{2026-10-16}: canonical status \texttt{partially\_solved}, evidence label \texttt{published}; parts (a) and (d) are described as archived solved parts, with the remaining parts in the main section. This is supplied dated curation, not an independently verified literature assessment.
\end{itemize}
\end{document}
